% Part: computability
% Chapter: computability-theory
% Section: ce-sets

\documentclass[../../../include/open-logic-section]{subfiles}

\begin{document}

\olfileid{cmp}{thy}{ces}
\olsection{গণনসাধ্যভাবে তালিকায়নযোগ্য সেট}

\begin{defn}
কোনো সেট শূন্য হলে অথবা কোনো গণনসাধ্য অপেক্ষকের মানপরিসর হলে সেটটিকে
\emph{গণনসাধ্যভাবে তালিকায়নযোগ্য} বলা হয়।
\end{defn}

\begin{history}
গণনসাধ্যভাবে তালিকায়নযোগ্য সেটকে \emph{পুনরাবৃত্তভাবে তালিকায়নযোগ্য}
সেটও বলা হয়। এটিই আদি পরিভাষা; আজ দুটি নামের সঙ্গে “c.e.” ও “r.e.”
সংক্ষিপ্ত রূপও প্রচলিত।
\end{history}

\begin{explain}
সংজ্ঞাটির অর্থ কী এবং পরিভাষাটি কেন যথাযথ, তা ভাবা উচিত। ধারণাটি হলো,
$S$ যদি গণনসাধ্য অপেক্ষক~$f$-এর মানপরিসর হয়, তবে
\[
S = \{ f(0), f(1), f(2), \dots \},
\]
তাই $f$-কে~$S$-এর উপাদানগুলি “তালিকায়িত” করছে বলে ধরা যায়। লক্ষ করুন,
সংজ্ঞা অনুসারে $f$-কে ক্রমবর্ধমান অপেক্ষক হতে হয় না, অর্থাৎ তালিকায়নটি
ক্রমবর্ধমান ক্রমে হওয়া আবশ্যক নয়। এমনকি $f$-কে একৈকও হতে হয় না;
অর্থাৎ~$S$-এর তালিকায়ন $f(0)$, $f(1)$, $f(2)$, \dots{}-এ পুনরাবৃত্তির
অনুমতি আছে। যেমন, ধ্রুবক অপেক্ষক $f(x) = 0$ সেট $\{ 0 \}$-কে তালিকায়িত
করে।
\end{explain}

প্রত্যেক গণনসাধ্য সেট গণনসাধ্যভাবে তালিকায়নযোগ্য। এটি দেখতে ধরা যাক,
$S$~গণনসাধ্য। $S$ শূন্য হলে সংজ্ঞা অনুসারেই সেটি গণনসাধ্যভাবে
তালিকায়নযোগ্য। অন্যথায়, $a$-কে~$S$-এর যেকোনো একটি উপাদান ধরি। $f$-কে
এভাবে সংজ্ঞায়িত করি:
\[
f(x) =
\begin{cases}
x & \text{যদি $\Char{S}(x) = 1$} \\
a & \text{অন্যথায়।}
\end{cases}
\]
তাহলে $f$ একটি গণনসাধ্য অপেক্ষক এবং~$S$ হলো~$f$-এর মানপরিসর।

\end{document}
